\subsection{The full intrinsic cubic class-two section set}

\begin{statement}{E.37}{Proved}{A complete local matrix criterion and every nonempty intrinsic fiber}
Use the same curve, tangent and group $N$ as in E.36, and put
\[
 V=N^{\rm ab},\quad M=\bigwedge^2V,\quad
 K=\mathbb Q(\zeta),\quad \lambda=1-\zeta .
\]
The following data give the entire image and every fiber of
\[
 \operatorname{Sect}(N/\gamma_3N\rtimes G_{\mathbb Q})
          /\text{geometric conjugacy}
       \longrightarrow H^1(\mathbb Q,V).
 \tag{E37.1}
\]
All coordinates in the criterion belong to the full rank-ten
module $V$.

Write the eight cusps other than $O=C_0$ in the order
\[
 \mathcal D=(A_0,A_1,A_2,B_0,B_1,B_2,C_1,C_2),
 \quad A_i=[0:-\zeta^i:1],\ B_i=[-\zeta^i:0:1],\
 C_i=[1:-\zeta^i:0].
\]
In the basis $(P,Q)=(A_0,A_2)$ of $E[3]$ their coordinates
are the rows of
\[
 \begin{pmatrix}1&0\\2&2\\0&1\\2&0\\1&1\\0&2\\2&1\\1&2\end{pmatrix}.
 \tag{E37.2}
\]
Let $e_d$ be the positive peripheral cocharacters of the
split torus $T_D$ over $K$, and write
$V=\mathbb Z_3f_1\oplus\mathbb Z_3f_2\oplus
\bigoplus_{d\in\mathcal D}\mathbb Z_3e_d$.
Set
\[
 R=\begin{pmatrix}0&-1\\1&-1\end{pmatrix},\quad
 C_E=\begin{pmatrix}1&-1\\0&-1\end{pmatrix},\quad
 U_d=(v_d,-u_d),\qquad (u_d,v_d)\text{ from (E37.2)}.
 \tag{E37.3}
\]
These are integral coordinates in which the unique homomorphism
$j:E\to J_{E,D}$ with $p\circ j=[3]$ has Tate matrix
$j_*=\binom{3I_2}{U}$.
For a prime $\mathfrak q\nmid3$ of $K$, let $q=N\mathfrak q$
and let $\psi_{\mathfrak q}=a+b\zeta$ be the elliptic
Frobenius scalar, equivalently the Jacobi sum of D.6.
It is the primary generator of $\mathfrak q$, congruent to
one modulo three; at an inert rational prime $p$ it is $-p$.
The actual Frobenius matrix on $V$ is
\[
 A_{\mathfrak q}=
 \begin{pmatrix}
  A_E&0\\ U(A_E-qI_2)/3&qI_8
 \end{pmatrix},\qquad A_E=aI_2+bR .
 \tag{E37.4}
\]
Write $A(\psi,q)$ for this displayed block matrix.
Let $B_{\mathfrak q}=\bigwedge^2A_{\mathfrak q}$.
The matrix of complex conjugation on $V$ is the same block
formula with $A_E=C_E$ and $qI_8$ replaced by $C_T$, where
$C_T$ is minus the permutation interchanging the subscripts
one and two in each cusp orbit.

Put
\[
 w=(3f_1+U_1)\wedge(3f_2+U_2),\quad
 L=\mathbb Z_3w\simeq\mathbb Z_3(1),\quad
 M_0=\ker(M\to\bigwedge^2T_3E).
 \tag{E37.5}
\]
The coefficient of $w$ at $e_{A_0}\wedge e_{A_2}$ is one.
Use as a basis of $M/L$ the other $44$ basic wedges,
and as a basis of $M_0$ the $44$ basic wedges other than
$f_1\wedge f_2$. The map
\[
 d_{\mathfrak q}:M/L\longrightarrow M_0,
 \qquad \overline m\longmapsto(q^{-1}B_{\mathfrak q}-1)m
 \tag{E37.6}
\]
is therefore a specified $44$ by $44$ matrix over $\mathbb Z_3$.
It is invertible over $\mathbb Q_3$.

Here is the arithmetic input for its right-hand side.
For each $d\ne O$, define $z_d\in J_{E,D}(K)$ to be
$\mathcal O(d-O)$, framed by $s_d^{-1}$ at $d$, by $s_O$
at $O$, and by one at every other cusp. The unit tangent
parameters are $s_{A_i}=x/z$, $s_{B_i}=y/z$,
and $s_{C_i}=z/x$.
Their third multiples are completely specified by
\[
 3z_d=(t_{d,r})_{r\in\mathcal D},\qquad
 t_{d,r}=\frac{\operatorname{lead}_r(f_d)}
                     {\operatorname{lead}_O(f_d)},\quad
 f_d=\frac{L_d}{x+y},
 \tag{E37.7}
\]
where
$L_{A_i}=y+\zeta^iz$, $L_{B_i}=x+\zeta^iz$,
$L_{C_i}=x+\zeta^{2i}y$.
At $d$ and $O$, the leading values mean respectively
$f_d/s_d^3$ and $f_ds_O^3$; elsewhere they mean ordinary
evaluation. Explicitly these exceptional values are
\[
 \begin{array}{c|cc}
 d&\operatorname{lead}_d(f_d)&\operatorname{lead}_O(f_d)\\ \hline
 A_i&1/3&3\\ B_i&1/3&-3\\
 C_i\ (i=1,2)&-1/(3(1-\zeta^i))&-3(1-\zeta^{2i}).
 \end{array}
 \tag{E37.8}
\]
Thus every entry in (E37.7) is an explicitly given unit
away from $\lambda$; the program also lists all $64$
entries as powers of $-\zeta^2$ and $\lambda$.

Represent $\alpha\in H^1(\mathbb Q,V)$ by a continuous
cocycle and restrict to $K$. At each $\mathfrak q\nmid3$
choose arithmetic Frobenius $F_{\mathfrak q}$ and a tame
pro-three inertia generator $\tau_{\mathfrak q}$ with
$F\tau F^{-1}=\tau^q$. Put
\[
 a_{\mathfrak q}=\alpha(F_{\mathfrak q}),\qquad
 v_{\mathfrak q}=\alpha(\tau_{\mathfrak q})
             =\sum_{d\in\mathcal D}b_{\mathfrak q,d}e_d .
 \tag{E37.9}
\]
The assertion that $v_{\mathfrak q}$ lies in this toric
submodule is part of the theorem. Let $k_{\mathfrak q,d}$
be any integral representative of the unramified Kummer
class of $z_d$ in
$V/(A_{\mathfrak q}-1)V$.
It is computable from the framed point $\overline z_d$
of $J_{E,D}(\mathbb F_q)$: take its successive
$3^m$-division points and apply $F_{\mathfrak q}-1$.
Define the following vector in the stated basis of $M_0$:
\[
 D_{\mathfrak q}
   =a_{\mathfrak q}\wedge v_{\mathfrak q}
      -\sum_{d\in\mathcal D}b_{\mathfrak q,d}
                         k_{\mathfrak q,d}\wedge e_d,
 \qquad
 c_{\mathfrak q}=-d_{\mathfrak q}^{-1}D_{\mathfrak q}.
 \tag{E37.10}
\]
Then the fiber of (E37.1) over $\alpha$ is nonempty
\emph{exactly} when
\[
 c_{\mathfrak q}\in\mathbb Z_3^{44}	ext{ for every }\mathfrak q\nmid3,
 \quad\text{and for every }m,\quad
 c_{\mathfrak q}\in3^m\mathbb Z_3^{44}
                   \text{ for all but finitely many }\mathfrak q.
 \tag{E37.11}
\]
All choices of the displayed Kummer representatives give
the same answer. The convergence requirement is included
in the criterion.

Every nonempty fiber is a torsor under the explicitly
evaluated group
\[
 \mathcal R_{\{p\ne3\}}(\mathbb Z_3)
       \oplus\mathbb Z_3^{24}
       \oplus(\mathbb Z/3)^9\oplus(\mathbb Z/9)^5 .
 \tag{E37.12}
\]
Here the restricted prime coordinates act by the actual
central cocycles $\sum_{p\ne3}a_p\kappa_3(p)w$.
For the remaining factor, put
$A_\tau=A(1+3\zeta,7)$,
$B_\tau=\bigwedge^2A_\tau$, and
$B_c=\bigwedge^2C_V$.
Its exact presentation, and hence a specified coordinate
map, is
\[
 H^1(G_{\mathbb Q,\{3,\infty\}},M)
   =M/(B_\tau-1)M^{c}
   =\operatorname{coker}\bigl((B_\tau-1)(1+B_c):M\to M\bigr).
 \tag{E37.13}
\]
The supplied exact Smith matrices have seven unit factors,
nine factors three, five factors nine, and $24$ zero
factors. Multiplication by the supplied left Smith matrix
gives the coordinates in (E37.12).

\emph{Manuscript: \S51.16, folios 733v--736v, scans 172--178,
(371)--(379). Equations (E37.1) and (E36.2) specify the
full intrinsic class-two group, including the $44$
directions discarded by the cuspidally central quotient.}
\end{statement}
\begin{proof}
\emph{The integral Tate lattice and its arithmetic action.}
The nine cusps are precisely $E[3]$ with origin $O$.
The tangent-line divisors prove that each is killed by
three, and they are nine distinct points. The collinear
cusp triples and the automorphism $z\mapsto\zeta z$
give (E37.2). On the Weierstrass model of E.24,
$P=(3,-5)$ and $Q=(3\zeta,-5)$.
The tangent functions $y+3x-4$ and $y+3\zeta^2x-4$
have divisors $3P-3O$ and $3Q-3O$, with equal leading
coefficients at $O$. The normalized Weil formula gives
\[
 e_3(P,Q)=-\frac{9(\zeta-1)}{9(\zeta^2-1)}=\zeta .
 \tag{E37.14}
\]
The sign in this formula is included in
\cite[Corollary 24.27]{E_SutherlandWeil}.

The framed-line-bundle description of the generalized
Jacobian identifies its extension characters with
$\mathcal O(d-O)$. Pulling back by $[3]$ trivializes
these characters, and gives the unique homomorphism $j$
with $pj=[3]$; uniqueness follows from
$\operatorname{Hom}(E,T_D)=0$.
For precision about its sign, if
$\operatorname{div}(h_d)=[3]^*(d-O)$, the vector frame
of the trivialized pullback is $1/h_d$. Translation by
$t\in E[3]$ multiplies this frame by
$h_d(X)/h_d(X+t)=e_3(d,t)^{-1}$.
Here $e_3$ is the divisor pairing used in (E37.14).
The equivalence with this translation formula, and
the contrary argument ordering in one presentation
of the pairing, are proved and discussed in
\cite[\S3, Remark 9]{E_EngePairings}.
Thus the character at $d$ of $j(t)$ is the inverse
divisor pairing. In the basis (E37.14), this gives
exactly $U_d=(v_d,-u_d)$.

Equivalently, the actual semiabelian variety is the
quotient of $T_D\times E$ by the graph
$\{(j(t)^{-1},t):t\in E[3]\}$, with projection to $E$
given by multiplication by three on the second factor.
The pairing computation proves that its extension
characters are the specified framed line bundles, so
this is an isomorphism of the given extensions.
Its Tate lattice is generated by the eight torus
generators and two lifts $f_1,f_2$ with
$j_*(t_i)=3f_i+U_i$. This proves (E37.3).
Nondegeneracy of (E37.14) shows that $j$ has trivial
kernel. Its image is an elliptic subvariety, and the
quotient is a torus; hence its Tate sublattice is
primitive. The wedge $w$ is also primitive, as witnessed
by its coefficient one at the indicated toric wedge.

The $\mathbb Z_3[\zeta]$-module $T_3E$ is torsion free
of rank one over this discrete valuation ring.
All its three-torsion is rational over $K$, so the
image of $G_K$ lies in $1+3\mathbb Z_3[\zeta]$.
Its determinant is the cyclotomic character by the
Weil pairing. The Jacobi-sum comparison D.6 gives the
scalar $\psi_{\mathfrak q}$ used in (E37.4).
Its primary congruence can also be checked directly:
the transformations $t\mapsto1-t$ and $t\mapsto1/t$
preserve the cubic character of $t(1-t)$.
Their orbits have sizes six, three, or two; the
size-two orbit satisfies $t^2-t+1=0$ and has character
one. Thus the character counts are $2,0,0$ modulo
three, and the negative character sum is one modulo
three. This works also in characteristic two.
Its reduction at $\mathfrak q$ vanishes, because the
corresponding power sum has degree
$2(q-1)/3<q-1$; its norm is $q$ by D.6.
These facts give the primary generator at split primes;
at inert primes the trace-fiber calculation in D.7
gives $-p$. Only one Eisenstein unit is one modulo
three, proving uniqueness.

Equation $A_Vj_*=j_*A_E$, together with the actual
torus action $qI_8$, now forces (E37.4).
To choose the stated complex-conjugation basis, first
lift the rational point $P$ to an
$\mathbb Z_3[\zeta]$-basis $t$ of $T_3E$.
It is a basis because $(\zeta-1)P\ne0$.
Write $c(t)=\epsilon t$. Then
$\epsilon\overline\epsilon=1$ and $\epsilon\equiv1$
modulo three. Its unique square root $v$ in
$1+3\mathbb Z_3[\zeta]$ satisfies $\overline v=v^{-1}$,
so $vt$ is fixed by $c$ and still lifts $P$.
Thus $c$ is conjugation on the Eisenstein coefficients,
with matrix $C_E$. The toric action is $C_T$, and
equivariance of $j_*$ gives the claimed remaining block.
These arguments determine all the displayed arithmetic
matrices on the integral lattice.

The rational decomposition induced by the isogeny
$T_D\times E\to J_{E,D}$ has Frobenius eigenvalues
$q$ eight times and the two elliptic eigenvalues
$\alpha_{\mathfrak q},\beta_{\mathfrak q}$.
On $M$ the only eigenvalue equal to $q$ is their
product, on the line $L$. The other eigenvalues have
complex absolute values $q^{3/2}$ or $q^2$.
Consequently (E37.6) is injective on $M/L$ and
invertible over $\mathbb Q_3$ onto $M_0$.
The index $[M:(L+M_0)]$ is nine, since the projection
of $w$ to $\bigwedge^2T_3E$ is nine. This explains why
the integral domain in (E37.6) must be $M/L$.
For example its determinant, up to a unit sign, is
\[
 \det d_{\mathfrak q}
    =\#E(\mathbb F_q)^8(q-1)^{28}/9.
 \tag{E37.15}
\]

\emph{Tangential paths and the complete peripheral term.}
Use a geometric path $p_d$ from the base tangent to
the chosen tangent at $d$. Its abelian path cocycle is
$k_d(g)=p_d^{-1}g(p_d)$ in $V$.
It is the Kummer class of the framed point $z_d$ above.
One can see the comparison on the actual covers as
follows. Choose an ordinary geometric base point and
the Abel--Jacobi map to the generalized Jacobian.
Pullbacks of its multiplication-by-$3^m$ covers are
the maximal abelian exponent-$3^m$ covers of the curve.
At a chosen tangent parameter, separate the power of
that parameter from the unit part of the local
framing. The coefficientwise Galois section fixes all
chosen powers of the parameter. Roots of the remaining
power series with constant term one are uniquely
normalized. Therefore the comparison of the two
tangential fibers is exactly the torsor of divisions
of the ratio of their constant vector frames.
There is an explicit local sign in that ratio. In
the family $\mathcal O(u_d-u_O)$, the canonical frame
at $d$ is $-t_d$ times $(s_d-t_d)^{-1}$, and at $O$
it is $(-t_O)^{-1}$ times $(s_O-t_O)$. Removing the
tangent powers therefore gives frames $-s_d^{-1}$
and $-s_O$, with one elsewhere. Their ratio to the
frames defining $z_d$ is the torus point $\epsilon_d$
with value $-1$ at $d$ and $O$ and one elsewhere,
modulo common scaling. This point has order two.
Its compatible $3^m$-division points are
$\epsilon_d$ itself, all rational, so its pro-three
Kummer torsor has the specified canonical trivialization.
Tensoring by that trivialization constructs an
isomorphism of the path torsor with the division
torsor of $z_d$ at every level, compatible with
every transition map.
It proves the integral Kummer comparison, including
the section origin; the path-cocycle convention is
the one in \cite[(12.4)--(12.10)]{E_WickelgrenThree}.

Multiplication by $f_d$ trivializes the third tensor
power of this framed line bundle. Evaluating the
frames gives (E37.7). Expanding the appropriate branch
$(1+s^3)^{1/3}$ gives (E37.8) directly. These formulas
prove that the points and their division torsors are
unramified away from $\lambda$; the cubic and its
distinct cusp divisor have good reduction there.
Over $\mathbb F_q$, the Kummer construction identifies
the three-primary part of $J_{E,D}(\mathbb F_q)$ with
$V/(A_{\mathfrak q}-1)V$. Explicitly a compatible
division sequence $y_m$ of a point gives
$((F-1)y_m)_m$. Changing the divisions changes this by
$(A_{\mathfrak q}-1)V$. Conversely $F-1$ is surjective
on the geometric semiabelian group, so the construction
gives every such class. It is therefore a finite-field
division computation of each required class, and does
not require any unknown class-two action.

Choose lifts of $f_1,f_2$ together with the eight
peripheral generators; they are a free pro-three basis
by the Frattini criterion. Since two is a unit in
$\mathbb Z_3$, the full class-two group in this basis
has the Baker--Campbell--Hausdorff law
\[
 (v,m)(v',m')=(v+v',m+m'+\tfrac12v\wedge v').
 \tag{E37.16}
\]
Write its arithmetic action as
$g(v,m)=(A_gv,B_gm+\eta_g(v))$.
Let $\delta_d$ be the local positive inertia loop.
Since the based peripheral loop is
$p_d^{-1}\delta_dp_d$, its Galois transform is
\[
 k_d(g)^{-1}(p_d^{-1}\delta_dp_d)^{\chi(g)}k_d(g).
\]
Thus, for $g\in G_K$, the exact term on the toric
submodule is
\[
 \eta_g(e_d)=-\chi(g)k_d(g)\wedge e_d.
 \tag{E37.17}
\]
This retains all eight cusp-path classes, with their
elliptic as well as toric components.

\emph{Local obstruction and the inverse limit.}
At $\mathfrak q\nmid3$ the genuine based action on the
geometric pro-three group is unramified, as established
in E.34. A cocycle restricted to tame inertia is
therefore a homomorphism. Its value $v$ satisfies
$A_{\mathfrak q}v=qv$. Projecting to $T_3E$ gives zero,
because the elliptic Frobenius has no eigenvalue $q$.
The toric sublattice is primitive, so $v$ belongs to
it integrally. Wild inertia has no nonzero homomorphism
to this pro-three group. This proves (E37.9).

Lift the inertia generator in the pulled-back
extension by $(v,m)$ and Frobenius by $(a,n)$.
The relation $F\tau F^{-1}=\tau^q$, computed with
(E37.16), gives the exact central equation
\[
 (B_{\mathfrak q}-q)m
       =-\eta_F(v)-q\,a\wedge v .
 \tag{E37.18}
\]
The central Frobenius coordinate $n$ cancels.
Dividing by $q$ and substituting (E37.17) gives
(E37.10). Inertia itself has the canonical zero-central
splitting $\tau^u\mapsto(uv,0)$, because the wedge of
two multiples of $v$ is zero. Thus (E37.18) is the
actual localization representative, including its
choice of inertia splitting, of the global central
extension obstruction.

Here is the continuous localization argument, with
the required support condition. Put
$S=\operatorname{Spec}\mathcal O_K[1/3]$ and
$M_n=M/3^nM$. The proof of E.34 gives
$H^2(G_{K,S},\mathbb F_3)=0$. The mod-three action on
$M$ has three-group image, so each finite $M_n$ has
a filtration with trivial $\mathbb F_3$ factors.
It follows that $H^2(S,M_n)=0$; global
three-cohomological dimension gives $H^3(S,M_n)=0$.
The finite-coefficient $H^1(S,M_n)$ are finite.
Localization and tame purity therefore give, exactly
as in E.32,
\[
 H^2_{\rm cts}(K,M)
  \simeq\operatorname{coker}\left(
   \mathcal R_{\mathfrak q\nmid3}(M(-1))
   \xrightarrow{\ F_{\mathfrak q}-1\ }
   \mathcal R_{\mathfrak q\nmid3}(M(-1))\right).
 \tag{E37.19}
\]
Here $\mathcal R$ means sequences tending to zero
modulo every $3^n$. To justify this formula, first
take the direct limit over the finite sets of removed
primes in the finite-coefficient localization triangles.
Their local two-term complexes have terms
$\bigoplus_{\mathfrak q}M_n(-1)$ and transition maps
which are surjective. Their derived inverse limit
is consequently the literal two-term complex of
restricted sequences. The two unramified groups in
degrees two and three vanish also after derived
inverse limit, by the preceding vanishing and
finiteness. This proves (E37.19) without discarding
a derived-limit class. The comparison and localization
results used here are those already opened in E.32,
\cite{E_MilneADT,E_JannsenContinuous}.

Continuity of $\alpha$ implies $v_{\mathfrak q}\to0$:
at every finite level it ramifies at only finitely
many primes. All $a_{\mathfrak q}$ and
$k_{\mathfrak q,d}$ are integral, so (E37.10) is a
restricted sequence. Its values lie in $M_0$ because
every summand has a toric wedge factor.
The kernel of each local differential is exactly $L$,
and the fixed integral section $M/L\to M$ obtained
by deleting the specified pivot wedge lifts restricted
sequences to restricted sequences. Therefore the
support representative is in the image in (E37.19)
exactly when its unique quotient corrections satisfy
(E37.11).

Changing $k_{\mathfrak q,d}$ by
$(A_{\mathfrak q}-1)t_{\mathfrak q,d}$ changes
$D_{\mathfrak q}$ by the differential of the integral
vector $-\sum_db_{\mathfrak q,d}t_{\mathfrak q,d}\wedge e_d$.
These vectors form a restricted sequence since
$b_{\mathfrak q,d}\to0$. Changing the cocycle
representative of $\alpha$ similarly changes $D$ by
the differential of an integral restricted sequence.
Frobenius and tame-generator choices are the changes
of the displayed actual local extension, so give
the same localization class. This proves the asserted
independence of all representatives.
Finally restriction from $\mathbb Q$ to $K$ is injective
in the obstruction group: corestriction composed with
restriction is multiplication by two, a unit on these
three-primary modules. Thus the criterion over $K$
is necessary and sufficient for a section over
$\mathbb Q$ itself.

\emph{All the fibers in explicit coordinates.}
E.36 proved $H^0(\mathbb Q,V)=0$. Hence the fiber over
a fixed abelian section class, once nonempty, is
exactly a torsor under $H^1(\mathbb Q,M)$, without a
further stabilizer quotient. We now compute this group.
At a rational prime $p\ne3$, the only Frobenius
eigenline of $M$ with eigenvalue $p$ is $L$.
The toric permutation eigenvalues do not alter this
conclusion: all the other complex absolute values
are $p^{3/2}$ or $p^2$. Its integral fixed lattice
after Tate twist is exactly $\mathbb Z_3w$ because
$w$ is primitive. Localization, with
$H^2(G_{\mathbb Q,\{3,\infty\}},M_n)=0$ proved by the
same filtration and quadratic descent, gives
\[
 0\to H^1(G_{\mathbb Q,\{3,\infty\}},M)
 \to H^1(\mathbb Q,M)
 \to\mathcal R_{\{p\ne3\}}(\mathbb Z_3)\to0.
 \tag{E37.20}
\]
The map $(a_p)\mapsto\sum_{p\ne3}a_p\kappa_3(p)w$
is an actual splitting: the sum converges uniformly
as a continuous cocycle and has precisely those
inertia residues.

It remains to evaluate the unramified term, not just
to name it. Let $P_S=G_{K,S}^{(3)}$, free pro-three of
rank two by E.34. The actual CM image is all
$1+3\mathbb Z_3[\zeta]$: Frobenius at the prime over
two is $-2$, and Frobenius at a prime over seven is
$1+3\zeta$ or its conjugate. Their logarithms divided
by three reduce to a basis modulo three, so they
topologically generate this group. Choose the prime
over seven giving $1+3\zeta$, and a Frobenius lift
$\tau$ in $P_S$. Then $\tau'=c\tau c$ has scalar
$1+3\zeta^2$. The logarithms of these latter two
elements also give a basis modulo three. Since
$P_S$ is free of rank two, the Frattini criterion
shows that $\tau,\tau'$ freely generate it.
Complex conjugation interchanges them exactly.

The based arithmetic action and all the cocycles
under consideration factor through
$P_S\rtimes\langle c\rangle$ as in E.34.
Every abelian cocycle may be normalized to vanish
on $c$, since two is invertible. In that normalization
it is determined freely by its value $x\in M$ on
$\tau$, and its value on $\tau'$ must be $B_cx$.
The coboundaries preserving the normalization are
exactly $(B_\tau-1)y$ with $y\in M^c$.
This proves the actual evaluation isomorphism
(E37.13). The map $1+B_c$ has image $M^c$ over
$\mathbb Z_3$, again because two is a unit.

The same description makes all abelian inputs of
(E37.10) concrete. Every proper elliptic cocycle is
unramified away from three: its inertia value would
otherwise give a Frobenius eigenvector of eigenvalue
$p$ in $T_3E$, which the good-reduction eigenvalues
exclude. Evaluation at $\tau$ gives
\[
 H^1(\mathbb Q,T_3E)
  =T_3E/(A_E(\tau)-1)(T_3E)^c
  =\mathbb Z_3t_1\oplus(\mathbb Z/3)t_2,
 \tag{E37.22}
\]
since $(T_3E)^c=\mathbb Z_3t_1$ and
$(A_E(\tau)-1)t_1=3t_2$.
For $b\in\mathbb Z_3$ and $r\in\{0,1,2\}$, the
unramified $V$-cocycle with value $bf_1+rf_2$ at
$\tau$, its conjugate value at $\tau'$, and zero
at $c$ is a specified lift of this elliptic class.
Subtracting that cocycle from an arbitrary $\alpha$
leaves a unique toric Kummer class. Uniqueness follows
from $H^0(\mathbb Q,T_3E)=0$, and the torus is
$\mathbb G_m^2\times(\operatorname{Res}_{K/\mathbb Q}
\mathbb G_m)^3$. Thus all inputs can equivalently
be supplied by $(b,r)$ and two completed rational
multiplicative classes and three completed
$K$-multiplicative classes. This is a bijection of
sets of inputs; it does not assert that this choice
of the representatives $r$ is a group splitting.
The Kummer classes have the prime and unit coordinates
computed in E.24, and the unramified cocycle is
determined on the two actual free generators just
specified. This provides evaluations of (E37.9)
without assigning a toric approximation to the
elliptic or peripheral-path coordinates.

The matrices (E37.2)--(E37.4) now compute the
cohomology by ordinary linear algebra over
$\mathbb Z_3$. The program performs exact rational
elementary operations whose denominators are prime
to three and saves both transformation matrices.
Their product with $(B_\tau-1)(1+B_c)$ is exactly
\[
 \operatorname{diag}(1^{\times7},3^{\times9},
                     9^{\times5},0^{\times24}).
 \tag{E37.21}
\]
This is an exact certificate for (E37.12), with the
coordinate map given by the saved left matrix.
For an independent check, the simultaneous invariants
of $M\otimes\mathbb Q_3/M$ under the matrices for
$-2$, $1+3\zeta$ and $c$ have elementary divisors
$3^{\times9},9^{\times5}$, exactly the torsion just
computed. Complex conjugation has minus-space
dimension $4\cdot6=24$, agreeing with the free rank.

The same program verifies all cusp frames, the Weil
sign, the primitive determinant line, and several
rank-$44$ determinant identities (E37.15). Formula
(E37.17), applied to the morphism of E.31 at the
prime above seven, gives the central residue two
computed there; this checks the peripheral minus
sign against an independent nonzero intrinsic
obstruction. The next maps for further lifting are
the actual surjections
$N/N''\to N/\gamma_3N$ and
$N/\gamma_3I_{\rm marked}\to N/\gamma_3N$,
with respective kernels
$\gamma_3N/N''$ and
$\gamma_3N/\gamma_3I_{\rm marked}$.
Those further section-lifting problems are not
asserted to vanish by the class-two calculation.
\end{proof}
